% 353_Koide_Amplitude_En_ch.tex
% Johann Pascher, 7 September 2026
\providecommand{\BM}{\textbf{[B]}}
\providecommand{\KM}{\textbf{[K]}}
\providecommand{\SM}{\textbf{[S]}}
\providecommand{\QM}{\textbf{[Q]}}

\phantomsection
\addcontentsline{toc}{chapter}{Koide Amplitude $d/c = \sqrt{2}$ from $T^4/Z_3$ Geometry}

\section*{Status Markers}
\begin{center}
\begin{tabular}{@{}lp{10.5cm}@{}}
\textbf{[K]} & \emph{Confirmed} --- derived from $\xi$, numerically verified. \\[4pt]
\textbf{[B]} & \emph{Proven} --- algebraically demonstrated. \\[4pt]
\textbf{[S]} & \emph{Speculative} --- open. \\[4pt]
\textbf{[Q]} & \emph{Source} --- external primary source, verified here. \\
\end{tabular}
\end{center}

\section*{Abstract}

The Koide condition $Q = 2/3$ is equivalent to the amplitude condition
$d/c = \sqrt{2}$ in the $Z_3$ circulant mass operator (Docs.~A110, 352).
This document derives $d/c = \sqrt{2}$ from the $T^4/Z_3$ geometry
without using Koide or the lepton masses as input.

The key is the trace bilinear form of GF$(27)^*$ (Doc.~348 Theorem~C):
the coupling matrix between the neutrino orbits ($f_1$) and the lepton orbits ($f_3$)
has normalised entries $|M^{(13)}_{ij}| = 1/\sqrt{2}$ at all non-zero positions.
This is a consequence of the $Z_3$ equal-distribution: each element couples to
exactly two of three partners, and the normalisation of the coupling strength
forces the weight $1/\sqrt{2}$ per connection.

The same condition appears in the mass operator as $d/c = \sqrt{2}$:
the circulant parameter $d$ is the coupling weight between the isotropic
part $c$ and the $Z_3$ interaction term, and the equal-distribution forces
$d/c = \sqrt{2}$.

Main findings: (1)~$d/c = \sqrt{2}$ from $Z_3$ equal-distribution \BM.
(2)~Connection to the GF$(27)^*$ coupling matrix $|M^{(13)}|_{\rm norm} = 1/\sqrt{2}$ \BM.
(3)~PDG satisfies $d/c = \sqrt{2}$ to $0.10\cdot\xi$ \KM.
(4)~Bare FFGFT residual $+17\cdot\xi$ explained by $\varepsilon_i$ ($101\,\%$) \KM.
(5)~Open R113-bridge (origin of falling $\varepsilon_i$ structure) \SM.

\section{Starting Point}

\subsection*{The Koide Amplitude Condition}

The $Z_3$ circulant mass operator (Doc.~A110) writes the square roots of the
lepton masses as:
\begin{equation}
  \sqrt{m_k} = c + d\cdot\cos\!\left(\theta + \frac{2\pi k}{3}\right),
  \quad k = 0,1,2.
  \label{eq:zirk}
\end{equation}
From $\sum_k \sqrt{m_k} = 3c$ and $\sum_k m_k = 3c^2 + \tfrac{3}{2}d^2$ it follows \BM:
\begin{equation}
  Q \equiv \frac{\sum m_k}{\bigl(\sum\sqrt{m_k}\bigr)^2}
  = \frac{1}{3} + \frac{d^2}{6c^2}.
  \label{eq:Q}
\end{equation}
The Koide condition $Q = 2/3$ is equivalent to \BM:
\begin{equation}
  \boxed{d = c\sqrt{2}.}
  \label{eq:bed}
\end{equation}
This condition does not contain $\theta$ --- $Q$ depends exclusively on
$d/c$ (Doc.~352 orthogonality theorem \BM).
The phase condition $\theta = 2/9$ is anchored in the geometry \KM;
the amplitude condition $d/c = \sqrt{2}$ was previously open \SM.

\subsection*{Known: $1/\sqrt{2}$ in the Galois Coupling Matrix}

Doc.~348 Theorem~C derives from the trace bilinear form
$\mathrm{Tr}_{\mathrm{GF}(27)/\mathrm{GF}(3)}(g^{k+j})$
the coupling matrix between neutrino orbits ($f_1$) and
lepton orbits ($f_3$) \BM:
\begin{equation}
  M^{(13)} = \begin{pmatrix} -1 & 0 & 1 \\ 1 & -1 & 0 \\ 0 & 1 & -1 \end{pmatrix},
  \qquad |M^{(13)}_{ij}|_{\rm norm} = \frac{1}{\sqrt{2}}
  \quad\text{at all non-zero entries.}
  \label{eq:M13}
\end{equation}
The weight $1/\sqrt{2}$ follows from the GF$(27)^*$ orbit structure.
The connection to the Koide amplitude condition is the subject of this document.

\section{Derivation of $d/c = \sqrt{2}$}

\subsection*{Step 1: The $Z_3$ Equal-Distribution Condition}

The $Z_3$ symmetry of the torus $T^4/Z_3$ arranges the three lepton generations
as equivalent elements of an orbit. A coupling between two $Z_3$ orbits
(e.g.\ neutrinos $f_1$ and charged leptons $f_3$) has the property, via the
trace bilinear form, that each element couples to exactly two of three elements
of the other orbit (the coupling matrix \eqref{eq:M13} has exactly one zero
entry per row).

The $\ell^2$ normalisation of the coupling strength per row gives:
\begin{equation}
  \sum_{j:\,M_{ij}\neq 0} |M_{ij}|^2 = 1
  \;\Longrightarrow\;
  2\cdot|M_{ij}|^2 = 1
  \;\Longrightarrow\;
  |M_{ij}| = \frac{1}{\sqrt{2}}.
  \label{eq:norm}
\end{equation}
This is the weight $1/\sqrt{2}$ from \eqref{eq:M13} \BM.

\subsection*{Step 2: Translation to the Circulant Mass Operator}

In the circulant mass operator \eqref{eq:zirk}, $c$ is the isotropic part
(mean of $\sqrt{m_k}$) and $d$ is the anisotropic part (amplitude of the
$Z_3$ modulation).

Parameter identification:
\begin{itemize}
  \item The isotropic part $c$ corresponds to the self-coupling of an orbit
    --- it sets the mass scale.
  \item The anisotropic part $d$ corresponds to the coupling between orbits
    --- it describes the interaction.
\end{itemize}

The weight of this coupling is fixed by the equal-distribution.
The coupling matrix \eqref{eq:M13} normalised to the isotropic scale:
every non-zero coupling has weight $1/\sqrt{2}$ \BM.

Therefore the ratio of the anisotropic to the isotropic part is:
\begin{equation}
  \frac{d}{c} = \frac{1}{|M_{ij}|_{\rm norm}^{-1}} = \frac{1}{1/\sqrt{2}} = \sqrt{2}.
  \label{eq:dc}
\end{equation}

This is $d/c = \sqrt{2}$ \BM.

\subsection*{Step 3: Geometric Reading}

The three lepton generations form a $Z_3$ orbit; the three neutrino generations
form a second. The trace bilinear form of GF$(27)$ describes how these orbits
couple to each other. The result is structurally forced:

\begin{enumerate}
  \item The $Z_3$ symmetry allows each element exactly $0$, $1$ or $2$
    non-zero couplings to the other orbit.
  \item The trace bilinear form selects exactly $2$ (Doc.~348 Theorem~C \BM).
  \item The $\ell^2$ normalisation forces the weight $1/\sqrt{2}$ per coupling.
  \item In the mass operator this weight appears as $d/c = \sqrt{2}$.
\end{enumerate}

The condition $d/c = \sqrt{2}$ is therefore not a fit to the measured value,
but a consequence of the $Z_3$ equal-distribution in the $T^4/Z_3$ geometry \BM.

\section{Numerical Verification}

\begin{center}
\begin{tabular}{@{}lllll@{}}
\toprule
Source & $d/c$ & $d/c - \sqrt{2}$ & Units $\xi$ & Status \\
\midrule
PDG \QM & $1.41420$ & $-1.0\times10^{-5}$ & $-0.098\cdot\xi$ & \KM \\
FFGFT bare & $1.41649$ & $+2.28\times10^{-3}$ & $+17.1\cdot\xi$ & \KM \\
FFGFT + $\varepsilon_i$ & $1.41422$ & $+8\times10^{-6}$ & $+0.06\cdot\xi$ & \KM \\
Target (geometry) & $\sqrt{2} = 1.41421$ & $0$ & $0$ & \BM \\
\bottomrule
\end{tabular}
\end{center}

PDG satisfies the condition to $0.10\cdot\xi$. The bare FFGFT residual
of $+17\cdot\xi$ is fully explained by the generation-dependent corrections
$\varepsilon_i$ (Doc.~352~\S9, $101\,\%$) \KM.
PDG masses: $m_e = 0.51100$~MeV, $m_\mu = 105.658$~MeV,
$m_\tau = 1776.86$~MeV \QM\,(PDG~2024).

\section{Status Overview}

\begin{center}
\begin{tabular}{@{}p{0.4cm}p{9.5cm}p{1.2cm}@{}}
\toprule
& Statement & Status \\
\midrule
1 & $Q = 2/3 \iff d/c = \sqrt{2}$ algebraically from $Z_3$ circulant & \BM \\
2 & $Q$ does not contain $\theta$; $d/c$ and $\theta$ are orthogonal parameters & \BM \\
3 & $d/c = \sqrt{2}$ from $Z_3$ equal-distribution of coupling to 2 of 3 & \BM \\
4 & Connection to GF$(27)^*$ coupling matrix $|M^{(13)}|_{\rm norm} = 1/\sqrt{2}$ & \BM \\
5 & PDG satisfies $d/c = \sqrt{2}$ to $0.10\cdot\xi$ & \KM \\
6 & Bare FFGFT residual $+17\cdot\xi$ explained by $\varepsilon_i$ ($101\,\%$) & \KM \\
7 & Dynamical derivation of the falling $\varepsilon_i$ structure & \SM \\
\bottomrule
\end{tabular}
\end{center}

Statement~3 closes the open bridge from Doc.~352 and register Doc.~190.
Status: raised from \SM\ to \BM.

\section{Note: Derivation Chain for Precise Lepton Masses}

\textbf{The lepton ladder is not in the chain for the precise mass derivation.}

The lepton ladder ($m_i = r_i\cdot\xi^{p_i}\cdot v$, Doc.~A100) yields masses
at the percent level (~1\,\%). It is a coarse framework, but not a necessary
link in the precise derivation. The [S] entry in Doc.~352~\S7
(``Koide is not contained in the FFGFT lepton ladder'') concerns the ladder
alone --- it is correct and remains [S]. It is, however, \emph{not} in the
chain for the precise masses.

\textbf{The chain for the precise lepton masses runs exclusively through the
$Z_3$ circulant:}

\begin{enumerate}
  \item $\theta = 2/9$ from the $Z_3$ circulant geometry \KM\ (Docs.~A110, 291)
  \item $d/c = \sqrt{2}$ from the $Z_3$ equal-distribution condition \BM\ (this doc.)
  \item Reference point $m_\mu/m_e$ \QM\ --- one declared empirical anchor
        (Doc.~A110, P42)
  \item $m_\tau = 1776.9690 \pm 0.0001$~MeV \KM\ (Docs.~A110, A210)
\end{enumerate}

This chain contains \emph{no} [S] entry. Doc.~A110 has no open points
in the circulant derivation. All [S] entries in Doc.~352 concern residual
analyses ($\varepsilon_i$ structure, R113-bridge) or the ladder--Koide
relation --- none of these is in the derivation chain for the precise masses.

\textbf{Conclusion:} The derivation chain for the precise lepton masses
is fully closed with the completion of $d/c = \sqrt{2}$ \BM\ (this document).
No [S] blocks it.

\section{Verification Script}

\begin{itemize}[leftmargin=2em,itemsep=2pt,topsep=2pt]
  \item Numerical verification of the $Z_3$ normalisation condition,
    coupling matrix entries \eqref{eq:M13}, $d/c$ calculation from PDG
    and bare FFGFT masses, $\varepsilon_i$ correction:
    \texttt{python/Dok353\_Skripte/pruef\_353\_koide\_amplitude.py}
\end{itemize}

\section{Open Points}

\begin{itemize}
  \item Why do $\varepsilon_i$ exhibit a falling generation-dependent structure
    ($\varepsilon_e > \varepsilon_\mu > \varepsilon_\tau$)?
    Attribution to QED extraction and SI projection is \SM\ (R113-bridge,
    Doc.~352~\S10).
  \item Full derivation of CKM/PMNS mixing angles from the Galois bilinear
    form (Doc.~348 Theorem~D) \SM.
  \item Are $c$ and $d$ individually fixed by the geometry, or only their
    ratio? \SM
\end{itemize}

\section{Use of AI Tools}

This document was produced with the assistance of AI-based tools.
The questions, postulates and all content decisions originate with the author;
the tools formulate, compute and generate verification scripts.
Responsibility for content and errors lies entirely with the author.
The detailed wording is in A010.
